% from ega4-11-fr.tex
\oldpage[IV-3]{133}

\begin{corollary}[11.3.2]
\label{IV.11.3.2-fr}
Soient $f:X\to Y$, $g:Y\to Z$ deux morphismes de présentation finie,
$\mathcal F$ un $\mathcal O_X$-Module quasi-cohérent de présentation finie.
Alors l'ensemble $U$ des $y\in Y$ tels que, pour toute générisation $y'$ de
$y$, $\mathcal F$ soit $(g\circ f)$-plat en tous les points de $f^{-1}(y')$
\textup{(i.e.} tels que $\mathcal F\otimes_Y\operatorname{Spec}(\mathcal O_{Y,y})$
soit plat relativement au morphisme
$X\times_Y\operatorname{Spec}(\mathcal O_{Y,y})\to Z$\textup{)} est ouvert dans $Y$.
\end{corollary}

Le même raisonnement que dans (11.1.5) montre que si $V$ est l'ensemble des
$x\in X$ où $\mathcal F$ est $(g\circ f)$-plat, $U$ est l'ensemble des $y\in Y$
dont toute générisation appartient à $E=Y\setmin f(X\setmin V)$. Comme $g\circ f$ est de
présentation finie, $V$ est ouvert dans $X$ en vertu de (11.3.1), donc
ind-constructible (1.9.6), et par suite $X\setmin V$ est pro-constructible~; mais comme
$f$ est quasi-compact, $f(X\setmin V)$ est pro-constructible dans $Y$ (1.9.5, (vii)),
et par suite $E$ est ind-constructible dans $Y$. Mais il résulte alors de
(1.10.1) que $U$ est l'intérieur de $E$, donc ouvert dans $Y$.

\begin{corollary}[11.3.3]
\label{IV.11.3.3-fr}
Soient $A$ un anneau, $B$ une $A$-algèbre de présentation finie, $C$ une
$B$-algèbre de présentation finie, $M$ un $C$-module de présentation finie.
Alors l'ensemble des $\mathfrak q\in\operatorname{Spec}(B)$ tels que
$M_\mathfrak q$ soit un $A$-module plat est ouvert dans
$\operatorname{Spec}(B)$.
\end{corollary}

\begin{proposition}[11.3.4]
\EGAReaderTarget{ega4.c04.prop.11-3-4}{IV.11.3.4-fr}
Soient $f:X\to S$ un morphisme localement de présentation finie, $\mathcal J$
un Idéal quasi-cohérent de type fini de $\mathcal O_X$, $Y$ le sous-préschéma
fermé de $X$ défini par $\mathcal J$, $\mathcal F$ un $\mathcal O_X$-Module de
présentation finie. On suppose que $\operatorname{gr}_{\mathcal J}^\bullet(\mathcal F)$
est $f$-plat. Dans ces conditions~:
\begin{enumerate}[label=\emph{(\roman*)}]
\item $\mathcal F$ est $f$-plat aux points de $Y$.
\item Si $\operatorname{gr}_{\mathcal J}^\bullet(\mathcal O_X)$ est $f$-plat,
$\operatorname{gr}_{\mathcal J}^\bullet(\mathcal O_X)$ est une
$(\mathcal O_X/\mathcal J)$-Algèbre de présentation finie et
$\operatorname{gr}_{\mathcal J}^\bullet(\mathcal F)$ est un
$(\operatorname{gr}_{\mathcal J}^\bullet(\mathcal O_X))$-Module de présentation
finie.
\item Supposons $\operatorname{gr}_{\mathcal J}^\bullet(\mathcal O_X)$ $f$-plat
(ce qui entraîne que $\mathcal O_X/\mathcal J$ est $f$-plat). Alors l'ensemble
$W$ des $y\in Y$ tels que $(\operatorname{gr}_{\mathcal J}^\bullet(\mathcal F))_y$
soit un $(\mathcal O_X/\mathcal J)_y$-module plat est ouvert dans $Y$.
\item Supposons $\operatorname{gr}_{\mathcal J}^\bullet(\mathcal O_X)$ $f$-plat.
Soit $S'\to S$ un morphisme quelconque~; soient $X'=X\times_SS'$,
$p:X'\to X$ la projection canonique, $Y'=p^{-1}(Y)=Y\times_SS'$ le
sous-préschéma fermé de $X'$, défini par
$\mathcal J'=p^*(\mathcal J)\mathcal O_{X'}$, $\mathcal F'=p^*(\mathcal F)
=\mathcal F\otimes_{\mathcal O_S}\mathcal O_{S'}$~; alors, si
$W'=p^{-1}(W)$, on a
$\operatorname{gr}_{\mathcal J'}^\bullet(\mathcal F')|W'
\cong p^*(\operatorname{gr}_{\mathcal J}^\bullet(\mathcal F)|W)$, et
$\operatorname{gr}_{\mathcal J'}^\bullet(\mathcal F')$ est un
$(\mathcal O_{X'}/\mathcal J')$-Module plat aux points de $W'$.
\end{enumerate}
\end{proposition}

\begin{proof}
Les questions étant locales sur $X$ et $S$, on peut supposer que
$S=\operatorname{Spec}(A)$ et $X=\operatorname{Spec}(B)$ sont affines, $B$ étant
une $A$-algèbre de présentation finie, et $\mathcal F=\widetilde{\mathbf M}$, où $M$ est
un $B$-module de présentation finie, $\mathcal J=\widetilde{\mathfrak J}$, où
$\mathfrak J$ est un idéal de type fini de $B$~; en vertu de (8.9.1) et
(8.5.11), il existe un sous-anneau noethérien $A_0$ de $A$, une $A_0$-algèbre
de type fini $B_0$ tels que $B=B_0\otimes_{A_0}A$, un idéal $\mathfrak J_0$ de
$B_0$ tel que $\mathfrak J=\mathfrak J_0B$, un $B_0$-module de type fini $M_0$
tel que $M=M_0\otimes_{A_0}A=M_0\otimes_{B_0}B$. En outre, $A$ est limite
inductive de ses sous-$A_0$-algèbres de type fini $A_\lambda$~; on pose
$B_\lambda=B_0\otimes_{A_0}A_\lambda$,
$M_\lambda=M_0\otimes_{A_0}A_\lambda=M_0\otimes_{B_0}B_\lambda$,
$\mathfrak J_\lambda=\mathfrak J_0B_\lambda$, de sorte que
$B=\varinjlim B_\lambda$, $M=\varinjlim M_\lambda$. Cela étant, par hypothèse
$\operatorname{gr}_{\mathfrak J}^\bullet(M)$ est un $A$-module plat~; donc il résulte
de (11.2.9) qu'il existe un $\lambda$ tel que
$\operatorname{gr}_{\mathfrak J_\lambda}^\bullet(M_\lambda)$ soit un
$A_\lambda$-module plat. Pour prouver que $\mathcal F$ est $f$-plat aux points
de $Y$, on peut donc se borner au cas où $S$ est \emph{noethérien}~; mais alors
$(\mathbf 0_{\mathrm I},6.1.2)$, appliqué par récurrence sur $n$, prouve que les
$\mathcal F/\mathcal J^{n+1}\mathcal F$ sont $f$-plats, et on conclut par
$(\mathbf 0_{\mathrm{III}},10.2.6)$.

